%!TEX root = ../hutbthesis_main.tex

% 设置英文摘要
\keywordsen{AirSim;  UAV swarm;  simulation platform;  performance optimization;  thread scheduling}
\begin{abstracten}

Aiming at the bottlenecks of sharp frame rate drop and surge delay in
large-scale UAV swarm simulation on Microsoft AirSim, this paper
conducts software-level optimization on UAV swarm simulation,
performance optimization and thread scheduling.  Inefficient thread
scheduling, frequent kernel-user mode switching, and unreasonable
spin-wait are identified as the main causes of performance degradation.
Five thread waiting models, including single-thread dispatcher, standard
blocking sleep, and cooperative yield, are designed and integrated into
the AirSim framework. The optimization effects are verified through
micro-benchmark and integrated simulation experiments from the
dimensions of timing accuracy, frame rate, and CPU usage. Results show
that the optimized scheme significantly improves the efficiency and
stability of AirSim for simulations with more than 20 UAVs, provides a
reusable tuning paradigm for high-concurrency real-time simulation,
reduces the cost of real-world testing for swarm algorithms, and
promotes the application of UAV swarms in reconnaissance, delivery, and
mapping.

\end{abstracten}
